% Part: normal-modal-logic
% Chapter: completeness
% Section: lindenbaums-lemma

\documentclass[../../../include/open-logic-section]{subfiles}

\begin{document}

\olfileid[id]{nml}{com}{lin}

\olsection{Lema Lindenbaum}

Lema Lindenbaum menetapkan bahwa setiap himpunan !!{formula}s yang
$\Sigma$-konsisten termuat dalam sekurang-kurangnya satu himpunan
\emph{lengkap} yang $\Sigma$-konsisten. Konstruksi model kanonik kita akan
menunjukkan bahwa bagi setiap himpunan lengkap dan
$\Sigma$-konsisten~$\Delta$, terdapat suatu dunia dalam model kanonik tempat
semua !!{formula}s dalam~$\Delta$---dan hanya formula-formula itu---benar.
Jadi, Lema Lindenbaum menjamin bahwa setiap
himpunan $\Sigma$-konsisten benar pada suatu dunia dalam model kanonik.

\begin{thm}[Lema Lindenbaum]\ollabel{thm:lindenbaum}
  Jika $\Gamma$ $\Sigma$-konsisten, maka terdapat suatu himpunan lengkap dan
  $\Sigma$-konsisten~$\Delta$ yang memperluas~$\Gamma$.
\end{thm}

\begin{proof}
Misalkan $!A_0$, $!A_1$, \dots{} merupakan suatu daftar menyeluruh semua
formula bahasa tersebut (pengulangan diperbolehkan). Sebagai contoh, mulailah
dengan mendaftarkan $\Obj p_0$, lalu pada setiap tahap $n \ge 1$ daftarkan
sejumlah berhingga formula dengan panjang~$n$ yang hanya menggunakan variabel
di antara $\Obj p_0$, \dots,~$\Obj p_n$. Kita mendefinisikan himpunan
!!{formula}s $\Delta_n$ melalui induksi pada~$n$, lalu menetapkan
$\Delta = \bigcup_n \Delta_n$. Pertama-tama, tetapkan
$\Delta_0 = \Gamma$. Dengan mengandaikan bahwa $\Delta_n$ telah didefinisikan,
kita mendefinisikan $\Delta_{n+1}$ sebagai berikut:
\[
\Delta_{n+1} =
\begin{cases}
  \Delta_n \cup \{!A_n\}, & \text{jika $\Delta_n \cup \{ !A_n\}$
    $\Sigma$-konsisten;} \\
  \Delta_n \cup \{ \lnot !A_n\}, & \text{jika tidak.}
\end{cases}
\]
Sekarang misalkan $\Delta = \bigcup_{n=0}^\infty \Delta_n$.

Kita harus menunjukkan bahwa definisi ini benar-benar menghasilkan himpunan
$\Delta$ dengan sifat-sifat yang diperlukan, yakni $\Gamma \subseteq \Delta$
dan $\Delta$ lengkap serta $\Sigma$-konsisten.

Jelas bahwa $\Gamma \subseteq \Delta$, sebab $\Delta_0 \subseteq \Delta$
berdasarkan konstruksi, dan $\Delta_0 = \Gamma$. Bahkan,
$\Delta_n \subseteq \Delta$ bagi semua~$n$, sebab $\Delta$ merupakan gabungan
semua~$\Delta_n$. (Karena pada setiap langkah konstruksi kita menambahkan
!!a{formula} ke himpunan yang telah dibangun, $\Delta_n \subseteq
\Delta_{n+1}$; jadi, karena $\subseteq$ transitif, $\Delta_n \subseteq
\Delta_{m}$ setiap kali $n \le m$.) Pada setiap tahap konstruksi, kita
menambahkan $!A_n$ atau $\lnot !A_n$, dan setiap !!{formula} muncul
(sekurang-kurangnya sekali) dalam daftar semua~$!A_n$. Jadi, untuk setiap
$!A$, berlaku $!A \in \Delta$ atau $\lnot !A \in \Delta$, sehingga $\Delta$
lengkap berdasarkan definisi.

Terakhir, kita harus menunjukkan bahwa $\Delta$ $\Sigma$-konsisten. Untuk itu,
kita menunjukkan bahwa (a) jika $\Delta$ tidak $\Sigma$-konsisten, maka suatu
$\Delta_n$ juga tidak $\Sigma$-konsisten, dan (b) semua $\Delta_n$
$\Sigma$-konsisten.

Andaikan $\Delta$ tidak $\Sigma$-konsisten. Maka $\Delta
\Proves[\Sigma] \lfalse$, yakni terdapat $!A_1$, \dots,~$!A_k \in
\Delta$ sedemikian sehingga $\Sigma \Proves !A_1 \lif (!A_2 \lif \cdots (!A_k
\lif \lfalse)\dots)$. Karena $\Delta = \bigcup_{n=0}^\infty \Delta_n$, setiap
$!A_i \in \Delta_{n_i}$ bagi suatu~$n_i$. Misalkan $n$ merupakan yang
terbesar di antaranya. Karena $n_i \le n$, $\Delta_{n_i} \subseteq \Delta_n$.
Jadi, semua $!A_i$ berada dalam suatu~$\Delta_n$. Hal ini berarti $\Delta_n
\Proves[\Sigma] \lfalse$, yakni $\Delta_n$ tidak $\Sigma$-konsisten.

Untuk menunjukkan bahwa setiap $\Delta_n$ $\Sigma$-konsisten, kita
menggunakan induksi sederhana pada~$n$. $\Delta_0 = \Gamma$, dan kita
mengasumsikan bahwa $\Gamma$ $\Sigma$-konsisten. Jadi, klaim tersebut berlaku
bagi $n = 0$. Sekarang andaikan klaim itu berlaku bagi $n$, yakni $\Delta_n$
$\Sigma$-konsisten. $\Delta_{n+1}$ adalah $\Delta_n \cup \{!A_n\}$ jika
himpunan tersebut $\Sigma$-konsisten; jika tidak, himpunan itu adalah
$\Delta_n \cup \{\lnot!A_n\}$. Dalam kasus pertama, jelas bahwa
$\Delta_{n+1}$ $\Sigma$-konsisten. Namun, berdasarkan
\olref[prf][con]{prop:consistencyfacts}\olref[prf][con]{prop:consistencyfacts-c},
$\Delta_n \cup \{!A_n\}$ atau $\Delta_n \cup \{\lnot!A_n\}$ konsisten,
sehingga dalam kasus lainnya pun $\Delta_{n+1}$ konsisten.
\end{proof}

\begin{cor}\ollabel{cor:provability-characterization}
  $\Gamma \Proves[\Sigma] !A$ jika dan hanya jika $!A \in \Delta$ bagi
  setiap himpunan lengkap dan $\Sigma$-konsisten~$\Delta$ yang memperluas
  $\Gamma$ (termasuk ketika $\Gamma = \emptyset$; dalam hal itu kita
  memperoleh pencirian lain bagi sistem modal $\Sigma$.)
\end{cor}

\begin{proof}
  Andaikan $\Gamma \Proves[\Sigma] !A$, dan misalkan $\Delta$ merupakan
  sebarang himpunan lengkap dan $\Sigma$-konsisten yang memperluas $\Gamma$.
  Jika $!A \notin \Delta$, maka berdasarkan kemaksimalan
  $\lnot!A \in \Delta$, sehingga $\Delta \Proves[\Sigma] !A$ (berdasarkan
  monotonisitas) dan $\Delta \Proves[\Sigma] \lnot!A$ (berdasarkan
  refleksivitas); dengan demikian, $\Delta$ tak konsisten. Sebaliknya, jika
  $\Gamma \Proves/[\Sigma] !A$, maka $\Gamma \cup \{ \lnot!A\}$
  $\Sigma$-konsisten, dan berdasarkan Lema Lindenbaum terdapat suatu himpunan
  lengkap dan konsisten~$\Delta$ yang memperluas
  $\Gamma \cup \{ \lnot!A \}$. Berdasarkan konsistensi, $!A
  \notin \Delta$.
\end{proof}

\end{document}
